% Generated by tools/edn_latex.py from reports/edn.md.
% Do not edit: change reports/edn.md and run `make edn`.
\ednentry{
  id = {EDN-45},
  title = {A \texorpdfstring{\texttt{req}}{req} marker is verification only where the specification says a test is the evidence},
  date = {2026-09-30},
  milestone = {M3: Quality Assurance},
  topic = {Testing Strategy},
  participants = {Lukas},
  decision = {The matrix reports four statuses, not two, and the specification's ``Verified by'' tag decides which one a \texttt{req} marker earns: a marker on an \textbf{[automated]} entry is \emph{verified by a test}, a marker on a \textbf{[manual]} entry is \emph{named by a test} and does not stand in for the drill its cell names. A marker on a \textbf{[manual]} entry whose cell names nothing does not satisfy the gate.},
  rationale = {\emph{Alternatives considered:}
\begin{itemize}[leftmargin=*, topsep=2pt]
\item \textbf{Option A (chosen): the specification's tag decides, and ``named by a test'' is a status of its own.}
\newline Pros: a fully checkable rule, no judgement in the tool; catches the real failure the review found, where NFR-06's determinism and MLflow provenance were reported as covered by a test whose whole body asserts that two files exist; closes the cheapest way to fake coverage, which was to tag an entry \textbf{[manual]}, leave its cell empty and put a marker on any test; makes the matrix state plainly that no status means a person agreed the evidence is enough.
\newline Cons: one more concept for a reader; a marker on a \textbf{[manual]} entry with named evidence still satisfies the gate, so the distinction changes the report and not the build in that case.
\item \textbf{Option B: keep one \texttt{covered} status, fix only the two wrong markers.}
\newline Pros: smallest diff; the two specific lies are gone.
\newline Cons: the mechanism that produced them stays, so the next marker on a \textbf{[manual]} entry reports as verified again; and the matrix keeps asserting something a tool cannot know.
\item \textbf{Option C: require the specification cell to name the test, and count a marker only when the named test exists.}
\newline Pros: the strongest claim of the three, and would also catch a marker on the wrong test.
\newline Cons: most cells name a category and not a test (``API test''), so it would need every cell rewritten before the API exists; ``does this cell name a test'' is a judgement a parser cannot make without a heuristic, and a heuristic in the middle of the traceability evidence is worse than a coarser rule that is exact.
\item \textbf{Option D: drop the automated/manual distinction from the gate and treat any marker as a claim needing human sign-off.}
\newline Pros: honest about the tool's limits.
\newline Cons: gives up the one thing the tool can check, and turns the gate back into a reminder.
\end{itemize}
\par
\emph{Why:} The matrix is what the report cites for NFR-07, so its worst failure mode is reporting a requirement as verified when nothing verified it, and that is exactly what happened: NFR-06 and NFR-01 were both green while their evidence was partly or wholly absent. A tool cannot read a test body and judge whether it verifies a criterion, so the honest move is to let the document that already says which evidence counts decide, and to say out loud in the generated table that a status is about the route being walked and not about a person agreeing it is enough. C would claim more than the documents can back today; B leaves the mechanism in place. The same reasoning is why the human review stays in NFR-07 rather than being replaced by this gate: NFR-01 still reads \emph{verified by a test} while two of its six success criteria have no thresholds, and only a person reading the table can notice that.},
  ai = {Alternative generation, Alternative assessment, Recommendation, Solution generation},
  response = {Accepted with modifications},
  assessment = {An adversarial AI review found the flaw and proved it by execution, naming the test whose body asserts that two files exist and the covering test whose own comment says four of the six criteria are still null. A second AI session generated A to D, recommended A over C on the grounds that C needs a heuristic, and implemented it. The modification is in the markers rather than the tool: AI's first instinct was to strip NFR-01's markers so the requirement reads as uncovered, which would have thrown away real evidence and turned the M3 gate red with nothing to fix; instead the marker moved onto the test that keeps SC-04 to SC-06 from ever counting as met, and the residual over-claim is stated in the pull request rather than hidden.},
  aievidence = {Claude Code adversarial review of PR \#51 on 2026-09-30, which reproduced the finding by running the generator and reading the covering tests; the follow-up session laid out A to D and verified each fix by re-running the reproduction.},
  evidence = {\href{https://github.com/mlops-2627q1-mds-upc/MLOPS_Recommenditos/pull/51}{PR \#51}; \texttt{Entry.\allowbreak{}status} and \texttt{Entry.\allowbreak{}has\_\allowbreak{}the\_\allowbreak{}promised\_\allowbreak{}evidence} in \texttt{tools/\allowbreak{}requirement\_\allowbreak{}matrix.\allowbreak{}py}; the status table in \texttt{CONTRIBUTING.\allowbreak{}md}; \href{https://github.com/mlops-2627q1-mds-upc/MLOPS_Recommenditos/issues/39}{issue \#39} for the SC-04 to SC-06 thresholds.},
}
